\chapter{Implementácia softvérovej časti}
\label{ImplementaciaSW}

\section{Firmvér ESP32}
\label{FirmverESP32}

Tu popíšem firmvér ktorý som napísal pre ESP32 uzly. Vysvetlím ako sa uzol pripája k Wi-Fi a k MQTT brokeru, ako prijíma a parsuje príkazy vo formáte JSON a ako na ich základe ovláda aktuátory. Priložím ukážku kľúčovej časti kódu.

\section{Backend Raspberry Pi}
\label{BackendRPI}

Tu popíšem backendovú aplikáciu bežiacu na Raspberry Pi. Vysvetlím štruktúru aplikácie -- triedy \texttt{MuseumController}, \texttt{ServiceContainer} a ostatné moduly -- a uvediem aké knižnice som použil. Priložím diagram tried.

\section{Stavový automat a~scene engine}
\label{StavovyAutomat}

Tu podrobne popíšem ako som implementoval stavový automat a scene engine. Vysvetlím ako sú reprezentované stavy a prechody, ako systém spracováva udalosti od návštevníka a ako scene engine sekvenčne vykonáva akcie definované v JSON konfigurácii.

\section{Webový dashboard}
\label{WebovyDashboard}

Tu popíšem webové rozhranie ktoré som vytvoril pre monitoring a manuálne ovládanie systému. Uvediem aký framework som použil a aké funkcie dashboard poskytuje. Priložím screenshot rozhrania.

\section{Konfigurácia scén -- JSON formát}
\label{KonfiguraciaScen}

Tu opíšem formát JSON súborov ktorými sa konfigurujú scény expozície. Na konkrétnom príklade vysvetlím všetky polia -- stavy, prechody, akcie -- a ukážem ako sa táto konfigurácia načítava a interpretuje systémom.
